\documentclass[12pt]{article}
\usepackage{amsmath, amssymb, amsthm, bm}
\usepackage[margin=1in]{geometry}
\usepackage[skip=0.6em]{parskip}

\newtheorem{lemc}{Lemma}[section]
\newtheorem{propc}[lemc]{Proposition}
\newtheorem{thmc}[lemc]{Theorem}
\theoremstyle{remark}
\newtheorem{remc}[lemc]{Remark}

\newcommand{\E}{\mathrm E}
\newcommand{\R}{\mathbf R}
\newcommand{\cC}{\mathcal C}
\newcommand{\cR}{\mathcal R}
\newcommand{\cF}{\mathcal F}
\newcommand{\wt}{\widetilde}

\begin{document}

\begin{center}
{\Large Toward Conjecture 12}\\[3pt]
{\small A best-effort attack: partial results, a barrier theorem, and a proof
plan :: August 29, 2026}
\end{center}

\medskip

This note does two things. Section~1 records my assessment of the summary
note for Fede. The rest is the honest state of my attack on Conjecture~12
($|\cC|\ge2$ implies every nice test is admissible): a master lemma that
unifies every admissibility proof we have (Section~2); a barrier theorem
showing that this method \emph{provably cannot} decide the remaining cases
(Section~3); two rigidity results that cut off the natural routes to a
counterexample (Sections~4 and~5); a reduction of the two-bad-facet case to
a single scalar inequality (Section~6); and the proof plan that all of this
points to (Section~7). I could not prove the conjecture. I believe the
material below is the reason no one-step proof exists, and it isolates the
one missing lemma. Compact proofs of the results stated without proof in the
summary note are collected in the appendix.

\section{Assessment of the summary note}

The note is accurate: I verified every stated result against the arguments
we have. Suggested repairs, all minor: the result numbering is
nonsequential (Lemma~1, Lemma~6, Definition~3, \dots) and should be
renumbered; ``or move involved'' should read ``or more involved''; in
Definition~3 the polyhedral criterion should note that irredundancy makes
$b_j=h_A(a_j)$, so the halfspaces are supporting; and the claim that the
redundant-only test is inadmissible for $\cC=\{1\}$ in every dimension
deserves its one-line justification --- the test depends only on
$(X_1,X_r)$, the projections of $\Lambda_0$ and $\Lambda_1$ fill the
bivariate null and alternative, and the bivariate dominator therefore lifts
(its level and power depend only on the $(X_1,X_r)$-marginal, which is
bivariate normal with correlation $\Omega_{1r}$, covered by the theorem for
every $\rho$). Two additions worth making: a sentence scoping the bivariate
necessity via $\Pi_\rho=\Omega_\rho\Lambda_1^*$, so the reader sees why
necessity is a $\Lambda_1^*\neq\{0\}$ phenomenon; and a description of the
open frontier --- regions with at least two essential infeasible facets,
with the canonical minimal examples (a rectangle with $|\cR|\ge2$; a
rectangle plus one binding infeasible cut; the tilted-gate family).

\section{The constraint set and the separating-score method}

Fix $|\cC|\ge2$ and a nice $A$, and let $\psi$ be a level-$\alpha$ test that
weakly dominates $\phi_A$; put $h:=\psi-\phi_A$ and
$M:=\E_0h^+=\E_0h^-$. Write $\wt A$ for the \emph{feasible hull} of $A$:
the intersection of its supporting halfspaces with normals in
$\cF(\Omega,\Lambda)$.

\begin{lemc}[the constraint set]\label{lem:constraints}
(i) $\{h>0\}\subseteq A$ and $\{h<0\}\subseteq A^c$;
(ii) $\E_0h=0$;
(iii) $q:=\E_0[Xh]=0$;
(iv) $\E_0[e^{\theta'X}h]\ge0$ for every $\theta\in\Omega^{-1}\Lambda_1$;
(v) $h^-=0$ a.s.\ on $\wt A^{\,c}$, so $h^-$ is supported on
$\wt A\setminus A$.
\end{lemc}

The proofs are those of the earlier notes (appendix). Every admissibility
result we have follows from Lemma~\ref{lem:constraints} through one scalar
comparison:

\begin{lemc}[separating score]\label{lem:score}
Suppose there exists
\begin{equation}\label{eq:score-class}
k(x)=c_0+\ell'x+\sum_{i=1}^m\nu_i\,e^{\theta_i'x},
\qquad
\nu_i\ge0,\ \ \theta_i\in\Omega^{-1}\Lambda_1,
\end{equation}
with $k\le0$ on $A$ and $k>0$ a.e.\ on $\wt A\setminus A$. Then $\phi_A$ is
$\alpha$-admissible.
\end{lemc}

\begin{proof}
By (ii)--(iv), $\E_0[kh]=c_0\E_0h+\ell'q+\sum_i\nu_i\E_0[e^{\theta_i'X}h]
\ge0$. By (i) and (v), $\E_0[kh]=\int_Akh^+ -\int_{\wt A\setminus A}kh^-
\le-\int kh^-<0$ whenever $M>0$. Hence $M=0$, $h=0$ $P_0$-a.s., and mutual
absolute continuity finishes the argument.
\end{proof}

Feasible support is the case $\wt A=A$ (no score needed); one essential bad
facet is the case $k=u'x-\kappa$, linear. The natural question is whether
Lemma~\ref{lem:score} can reach the frontier regions. It cannot.

\section{A barrier theorem}

Take the canonical frontier case: $\Omega=I_4$, $\cC=\{1,2\}$,
$\cR=\{3,4\}$, and the exact-size rectangle
$A=\bigcap_{j=1}^4\{x_j\le\kappa_j\}$, whose two bad facets are the
redundant faces $e_3,e_4$.

\begin{thmc}[no separating score exists]\label{thm:barrier}
For this $A$, no function $k$ of the form~\eqref{eq:score-class} satisfies
$k\le0$ on $A$ and $k>0$ on $\wt A\setminus A$.
\end{thmc}

\begin{proof}
First, each $\theta_i$ must be nonnegative: if $\theta_{ij}<0$, move from
any $y^0\in A$ along $-e_j$ (which stays in $A$); the $i$th term grows
exponentially while $|\ell'x|$ grows linearly and the other terms are
nonnegative, so $k>0$ somewhere on $A$. Being also in $\Lambda_1$, each
$\theta_i$ is feasible, and any nonzero feasible $\theta\ge0$ has
$\theta_1+\theta_2>0$ (if $\theta_1=\theta_2=0$, feasibility forces
$\theta_3=\theta_4=0$). Second, $k\le0$ on $A$ as $x_j\to-\infty$ forces
$\ell\ge0$.

Now consider the gap points $z_T=(-T,-T,\kappa_3+1,\kappa_4-1)$. They
violate the third face, satisfy the others, and lie in $\wt A$ for large
$T$: for feasible $\theta$,
$\theta'z_T-\theta'\kappa
=\theta_3-\theta_4(\cdot)-(T+\kappa_1)\theta_1-(T+\kappa_2)\theta_2
\le(\theta_1+\theta_2)(1-T-\min_j\kappa_j)\le0$
using $\theta_3\le\max(\theta_1,\theta_2)\le\theta_1+\theta_2$. Since every
$\theta_i$ has $\theta_{i1}+\theta_{i2}>0$, all exponential terms vanish as
$T\to\infty$, so $k(z_T)=c_0+\ell'z_T+o(1)$. Positivity for all $T$ forces
$\ell_1+\ell_2\le0$, hence $\ell_1=\ell_2=0$; the variant with
$x_4\to-\infty$ (still in the gap and in $\wt A$) forces $\ell_4=0$, and
the symmetric fourth-face lobe forces $\ell_3=0$. So $\ell=0$ and
$k=c_0+\sum_i\nu_ie^{\theta_i'x}$. If $c_0>0$, then $k(z_T)\to c_0$ is fine
but deep points of $A$ give $k\to c_0>0$, violating $k\le0$ on $A$. If
$c_0<0$, then $k(z_T)\to c_0<0$, violating gap positivity. If $c_0=0$, then
$k\ge0$ everywhere, and $k\le0$ on $A$ forces $\nu=0$, i.e.\ $k\equiv0$,
violating gap positivity. Contradiction in every case.
\end{proof}

The same argument runs for the three-dimensional tilted-gate family
$A(\lambda)=\{x_1\le\kappa,x_2\le\kappa,x_1+\lambda x_3\le b,
x_2+\lambda x_3\le b\}$, $\lambda>1$: each gate's ``dead lobe'' (both cores
escaping to $-\infty$ along the gate) forces $\ell$ proportional to that
gate's normal, and the two lobes together force $\ell=0$, after which the
constant-term trichotomy applies.

The interpretation is sharp. Every admissibility proof we possess
aggregates Lemma~\ref{lem:constraints} into one scalar comparison, and
Theorem~\ref{thm:barrier} says no such aggregation can decide the frontier.
The obstruction lives in the \emph{dead zones}: gap regions escaping to
$-\infty$ in all core coordinates, where every alternative tilt vanishes
and only the (sign-constrained) linear terms reach --- and two bad facets
make incompatible demands on the single linear term. A proof of
Conjecture~12 therefore needs either a new confinement principle that
empties the dead zones before scoring, or a genuinely non-scalar use of the
constraints; and a counterexample, if one exists, must place comparison
mass in the dead zones. The next two sections squeeze the second
possibility.

\section{No dominator hides in a dead slab}

\begin{propc}[slab rigidity]\label{prop:slab}
Let $\Omega=I_d$ and let $A$ be the exact-size rectangle (any $|\cC|\ge1$,
any $|\cR|$). Fix $T>0$ and suppose $h=1\{x_1\le-T\}\,g(x_2,\dots,x_d)$
satisfies the conditions of a dominance perturbation
(Lemma~\ref{lem:constraints}(i)--(ii) and $\E_\mu h\ge0$ for all
$\mu\in\Lambda_1$). Then $g=0$ a.s.
\end{propc}

\begin{proof}
By independence, $\E_\mu h=\Phi(-T-\mu_1)\,\E_\nu g$ with
$\nu=(\mu_2,\dots,\mu_d)$. As $\mu_1\to\infty$ the constraint
$\mu_r\le\max_c\mu_c$ frees $\nu$ completely, so $\E_\nu g\ge0$ for
\emph{every} $\nu\in\R^{d-1}$, and (ii) gives $\E_0g=0$. The sign structure
makes $g$ a dominance perturbation of the $(d-1)$-dimensional section
rectangle against the \emph{unrestricted} alternative. In the unrestricted
problem every direction is an alternative direction, so every closed convex
acceptance region is feasibly supported, and the confinement argument of
Lemma~\ref{lem:constraints}(v) applies to every supporting halfspace of the
section: $g^-=0$ a.s., and $\E_0g=0$ then gives $g=0$ a.s.
\end{proof}

So the dead zones, where the barrier lives, cannot host a self-contained
dominator: a slab of a $|\cC|\ge2$ problem faces an effectively
unrestricted alternative, which is the Birnbaum regime where convexity
already wins. Any dominator must couple dead zones to live zones.

\section{The natural coupling fails: localized gates}

There is exactly one natural family of couplings: borrow the $|\cC|=1$
gated dominator and arm it only when the data indicate the other core is
dead. In the three-dimensional two-core model with the equal-cutoff
rectangle, let $g_{cr}$ denote the calibrated gate swap in coordinates
$(c,r)$ (reject the diagonal corner $G_{cr}$, accept the strip $D_{cr}$,
with $P_0(G_{cr})=P_0(D_{cr})$), and define
\[
h_L
:=
1\{x_1\le-L\}\,g_{23}(x_2,x_3)
+
1\{x_2\le-L\}\,g_{13}(x_1,x_3)~.
\]
This satisfies the sign structure and $\E_0h_L=0$, and it gains where it was
designed to: numerically, at $\mu=(2,-5,1)$ (core~1 active, core~2 dead),
$\E_\mu h_L\simeq+0.108$. But it violates the first-moment identity:
$q_2=\Phi(-L)\,\E_0[X_2g_{23}]>0$, and Lemma~\ref{lem:constraints}(iii)
says a dominator must have $q_2=0$. The violation is realized at near-null
alternatives with a dead core: at $\mu=(t,-3,0)\in\Lambda_1$ the armed
$(2,3)$-gate loses a fixed amount, $\Phi(-L)\Delta_{23}(-3,0)<0$, while the
compensating $(1,3)$-gain vanishes linearly in $t$. Numerically with
$L=2$, $\alpha=0.05$: $\E_\mu h_L\simeq+0.0113$ at $t=0.5$, $+0.0017$ at
$t=0.1$, and $-2.8\times10^{-4}$ at $t=0.005$. Repairing this requires a
switch with $\E_0[\text{switch}]=0$, which the sign structure forbids (the
negative part of the switch times the negative part of the gate puts
$h>0$ outside $A$). I regard this as the strongest concrete evidence for
the conjecture: the only known source of power improvements, transplanted
in its most favorable form, is killed by the first-moment identity exactly
where that identity has bite.

\section{The two-facet case is one scalar inequality}

\begin{lemc}[absolute-value reduction]\label{lem:absred}
Let $A=\wt A\cap\{g_u\le0\}\cap\{g_w\le0\}$ be nice with two essential bad
facets, where $g_u(x)=u'x-\kappa_u$, $g_w(x)=w'x-\kappa_w$. If $M>0$, then
\[
\E_0\bigl[\,|g_u(X)-g_w(X)|\,h\bigr]<0~.
\]
Consequently, Conjecture~12 holds for all such $A$ if and only if every
dominance perturbation satisfies
$\E_0[\,|(u-w)'X-(\kappa_u-\kappa_w)|\,h]\ge0$.
\end{lemc}

\begin{proof}
Let $f=\max(g_u,g_w)$. Then $f\le0$ on $A$ and $f>0$ a.e.\ on
$\wt A\setminus A$, so as in Lemma~\ref{lem:score}, $M>0$ implies
$\E_0[fh]<0$. But $f=\tfrac12(g_u+g_w)+\tfrac12|g_u-g_w|$, and the affine
part integrates to zero against $h$ by (ii)--(iii).
\end{proof}

For $A(\lambda)$ the relevant direction is $u-w\propto(1,-1,0)$, and
\emph{both} $\pm(1,-1,0)$ lie in $\Lambda_1$, so two-sided exponential
information $\E_0[\cosh\{t(X_1-X_2)\}h]\ge0$ is available for every $t$.
The missing step is exactly that $|z|$ is not a mixture of
$\cosh(tz)$'s and affine functions --- every such mixture grows
superlinearly --- which is Theorem~\ref{thm:barrier} again in scalar form.
For the four-dimensional rectangle, $u-w=e_3-e_4$ and neither sign is in
$\Lambda_1$, so even this reduction has less to work with; the rectangle
with $|\cR|\ge2$ is, I believe, the hardest case.

\section{The proof plan}

The results above locate the single missing ingredient.

\emph{Step 1 (the key lemma: dead-zone confinement).} Prove that every
dominance perturbation has $h=0$ a.s.\ on the dead zones --- gap regions
along which all coordinates in some core-detecting set escape to
$-\infty$. Proposition~\ref{prop:slab} proves this for product-form $h$;
the open problem is the coupled case. The natural tool is a conditional
argument in the spirit of the bivariate rearrangement lemma: condition on
the deep core coordinates and use the fact that the conditional problem is
effectively unrestricted, where the section of a nice region is convex and
hence rigid. The technical difficulty is that dominance constrains only
the $x_1$-integral of the conditional perturbations, not each slice.

\emph{Step 2 (score on the truncated gap).} Granted Step~1, $h^-$ is
confined to the part of $\wt A\setminus A$ where some feasible tilt stays
active. There the barrier computation no longer applies, and finite
mixtures~\eqref{eq:score-class} plausibly separate: each lobe of the
truncated gap is exposed by tilts near the feasible boundary of its gate
normal. This should be checked first numerically (a small convex program:
minimize the gap infimum of $k$ minus the $A$-supremum over
$(c_0,\ell,\nu,\theta)$), then proved for the canonical regions.

\emph{Step 3 (the decisive experiment, in parallel).} The discretized
linear program from the earlier note, aimed at $A(2)$ and at the
four-dimensional rectangle, with the $1\leftrightarrow2$ symmetry imposed
and the alternative net including deep rays. A robustly positive value
yields the counterexample; a vanishing value yields, through the dual
solution, the weights of the separating construction that Step~2 needs ---
including any capacity multipliers (dual variables on $0\le\psi\le1$),
which are exactly the terms the score class~\eqref{eq:score-class} lacks
and which a Blyth-type argument would build implicitly.

\emph{Where I stand.} Sections~4 and~5 close the two natural routes to a
counterexample; Section~3 explains why no proof has appeared. My credence
in Conjecture~12 has gone up: the dead zones are the only place a
dominator could live, and they are inert both in isolation
(Proposition~\ref{prop:slab}) and under the one natural coupling
(Section~5). But the conjecture is genuinely open, and Step~1 is where I
would spend effort --- it is a self-contained lemma, its product-form case
is already proved, and it converts the barrier from an obstruction into a
two-line finish.

\appendix

\section{Proofs of the results stated in the summary note}

\emph{Lemma 1 (dual-cone dichotomy).} For $\cC=\{1\}$: normalizing
$\mu_1=1$, $v_1+\sum_r\mu_rv_r\ge0$ for all $\mu_r\le1$ forces $v_r\le0$
(let $\mu_r\to-\infty$) and then $v_1+\sum_rv_r\ge0$ (at $\mu_r=1$);
conversely these suffice, and the cone representation is a direct check.
For $|\cC|\ge2$: $e_c+se_{c'}\in\Lambda_1$ for all real $s$ kills the core
coordinates pairwise; means with one redundant coordinate at the core
maximum and the others at $-M$ kill the redundant coordinates.

\emph{Lemma 6 (first-order restriction).} $\mu=0\in\Lambda_0$ gives
$\E_0\psi\le\alpha=\E_0\phi$; dominance along $te_c\in\Lambda_1$ and
continuity of Gaussian power functions give the reverse, so $\E_0h=0$.
For $\mu\in\Lambda_1$, $t\mapsto\E_0[he^{t(\Omega^{-1}\mu)'X}]$ is
nonnegative for $t>0$ (dominance along $t\mu$, a ray in the cone
$\Lambda_1$), vanishes at $t=0$, and is differentiable there ($|h|\le1$,
Gaussian m.g.f.), so its right derivative $(\Omega^{-1}\mu)'q$ is
nonnegative; hence $\Omega^{-1}q\in\Lambda_1^*$.

\emph{Proposition 4 (feasible support suffices).} For nonzero
$a\in\cF(\Omega,\Lambda)$ with $h_A(a)<\infty$: $\Omega a\in\Lambda$ and
$a'\Omega a>0$ put $\Omega a\in\Lambda_1$, so
$\E_0[e^{ta'X}h]\ge0$ for all $t>0$; since $h^+$ lives on
$A\subseteq\{a'x\le h_A(a)\}$, letting $t\to\infty$ kills
$h^-$ above the supporting plane. Lindel\"of reduces the union of the
half-space complements to a countable one, so $h^-=0$ a.s.\ on $A^c$ when
$A$ is feasibly supported; $\E_0h=0$ then gives $h=0$ a.s.

\emph{Proposition 5.} If all finite normals are feasible, the feasible
supporting halfspaces include all supporting halfspaces, whose intersection
is $A$. Conversely with $|\cC|=1$, $\cF$ is a closed convex cone (the
defining constraint is linear), the barrier cone of $A$ lies in the closed
convex cone generated by the feasible normals used in the representation,
hence in $\cF$.

\emph{Proposition 7 (one arbitrary halfspace).} Confinement puts $h^-$ on
$\{u'x>\kappa\}\cap A_f$; $h^+$ lives on $A\subseteq\{u'x\le\kappa\}$; if
$M>0$ then $\E_0[(u'X-\kappa)h]<0$, contradicting $u'q=0$ from Lemma~6 and
the dichotomy.

\emph{Corollaries 8--9 and core-only rectangles.} Corollary~8: $A_f=\R^d$,
$u=a$. Corollary~9: core faces are feasible since
$(\Omega e_c)_r=\Omega_{rc}<1=(\Omega e_c)_c$; the single redundant face is
the extra halfspace. Core-only rectangles: all facets feasible by the same
computation, so Proposition~4 applies for any $|\cC|\ge1$.

\end{document}
